\bigskip
\centerline{\bf 3. Derivatives along Recurrent  Orbits}
\bigskip

\flushpar
Let $\rho_n=\min\{\frac{1}{\sigma_{n-1}},\frac{1}{\sigma_n}\}$.
In this section we will prove 

\proclaim{Proposition 3.1} There exist  $n_0\ge 1$, $\theta<1$ and $C>0$ with 
the following property. If $n\ge n_0$, 
$x\in U_{n+1}$ and $R_n^i(x)\notin U_{n+1}$ for $i\le s$
then
$$
|Df^T(f(x))|\ge C\cdot \rho_n \cdot \theta^{-(s-1)},
$$
where $R^s_n(x)=f^T(x)$.
\endproclaim

\flushpar
In [VT] a similar estimate in the case $s=1$ was obtained for circle 
homeomorphisms with irrational rotation number of bounded type.
The proof of Proposition 3.1 will use the following Lemmas 
and notation. Fix
$x\in U_{n+1}$ according to the Proposition, say $R_n^i(x)=f^{t_i}(x)$ with
$i\le s$.

\proclaim{Lemma 3.2} For each $i\le s$ there exists an interval $S_i\ni f(x)$
such that
$$
f^{t_i-1}:S_i\to U_n
$$ 
is monotone and onto.
\endproclaim

\demo{Proof} Lemma 2.2 applied to $U_{n+1}\in \Cal{U}_n$ states the existence 
of $S_1$. Assume that $S_i\ni x$
exists. Then $f^{t_i-1}: S_i\to U_n$ monotone and onto. Moreover 
$f^{t_i-1}(f(x))\in I_{i+1}\in \Cal{U}_n$. Because 
$f^{t_i-1}(f(x))\notin U_{n+1}$ we have that $I_{i+1}\ne U_{n+1}$ and 
$R_n:I_{i+1}\to U_n$ is monotone and onto. Now let 
$S_{i+1}=f^{-(t_i-1)}(I_{i+1})\cap S_i$.
\hfill\hfill\qed $\,\,$ (Lemma 3.2)
\enddemo

\flushpar
Observe that $f^{t_{i-1}-1}(S_i)=I_i\in \Cal{U}_n$.

\proclaim{Lemma 3.3} There exist $n_0\ge 1$ and $K<\infty$ with the 
following property. If the distortion of
$$
f^{t_i-1}:S_i\to U_n \text{ with } n\ge n_0
$$ 
is bigger than $K$ then
$$
I_i\subset (-\frac34 \cdot u_n, \frac34 \cdot u_n).
$$
\endproclaim

\demo{Proof} Lemma 2.2 states that $f^{t_1}:S_1\to U_n$ has a monotone 
extension up to $U_{n-1}$, the map is essentially linear for big enough $n$. 
Hence $i\ge 2$.
Consider the following decomposition
$$
f^{t_i-1}|S_i=f^{t_{i}-t_{i-1}-1}|f(I_i)\circ 
              f|I_i\circ
              f^{t_{i-1}-1}|S_i.
$$ 
The factor $f^{t_{i-1}-1}|S_i$ has a monotone extension up to 
$U_n$. In particular, for big enough $n$, it
is essentially linear. This is because the image $I_i$
has a small hyperbolic length within $U_n$ (Lemma 2.3). The factor
$f^{t_{i}-t_{i-1}-1}|f(I_i)$
has a monotone extension up to $U_{n-1}$ (Lemma 2.2), it is also essentially
linear. The distortion of $f^{t_i-1}|S_i$ can only be caused by the factor
$f|I_i$, $I_i$ has to be very close to $0$. There is some $n_0$
such that $I_i\subset (-\frac34 \cdot u_n, \frac34\cdot u_n)$, 
whenever $n\ge n_0$.
Here we used Lemma 2.3 which states that $I_i$ has also very small hyperbolic 
length within $U_n$.
\hfill\hfill\qed $\,\,$ (Lemma 3.3)
\enddemo

\proclaim{Lemma 3.4} For any $\theta<1$ there exist $n_0\ge 1$ and 
$C<\infty$ such that 
$$
|S_i|\le C\cdot u_{n+1}^\alpha \cdot \theta^{i-1},
$$
whenever $n\ge n_0$ and $i\ge 2$.
\endproclaim

\demo{Proof} Observe that $f(0)\in S_1\supset S_2\supset\dots\supset S_i$
and $S_i\subset f(U_{n+1})\subset S_1$, $i\ge 2$.
Let $L_j\subset S_1$ be the connected component of $S_1-S_j$ with
$L_j\subset f(U_{n+1})$, $2\le j\le i$.  For $n$ big enough we get 
from the proof of Lemma 3.2 and from Lemma 2.3 that the hyperbolic length of 
$S_j$ within $S_{j-1}$ is very small, $2\le j\le i$. It is easily seen that
this implies 
$$
\aligned
|S_i|&\le C\cdot \theta^{i-1}\cdot |L_i|\\
   &\le C\cdot \theta^{i-1}\cdot |f(U_{n+1})|\\
   &\le C\cdot \theta^{i-1}\cdot u_{n+1}^\alpha.
\endaligned
$$
\hfill\hfill\qed $\,\,$ (Lemma 3.4)
\enddemo

\proclaim{Lemma 3.5} There exist $n_0\ge 1$ and  $C>0$ such that the following 
holds for $n\ge n_0$. If $|R_n(U_{n+1})|=|f^{t_1}(U_{n+1})|\ge 
\frac{1}{10} u_n$ then
$$
|Df^{t_1-1}_{|S_1}|\ge C\cdot \frac{u_n}{u_{n+1}^\alpha}.
$$
If $|R_n(U_{n+1})|<\frac{1}{10} u_n$ then
$$
|Df^{t_1-1}_{|S_1}|\ge C\cdot \frac{u_{n-1}}{u_n^\alpha}.
$$
\endproclaim

\demo{Proof} Consider the map $f^{t_1-1}: S_1\to U_n$. From Lemma 2.2 we know 
that this map has a monotone extension up to $U_{n-1}$. The map is essentially
linear because $u_{n-1}>>u_n$ whenever $n$ is big enough. 
The derivative $|Df^{t_1-1}_{|S_1}|$ is
essentially constant and can be estimated by
$$
|Df^{t_1-1}_{|S_1}|\ge C\frac{|R_n(U_{n+1})|}{|f(U_{n+1})|}
                        \ge C\frac{u_n}{u_{n+1}^\alpha}.
$$
Here we used that $|R_n(U_{n+1})|\ge \frac{1}{10} u_n$.
Now consider the other case: 
$|R_n(U_{n+1})|< \frac{1}{10} u_n$. Let $K\supset f(U_{n+1})$
be the interval which is mapped monotonically onto $U_{n-1}$:
$f^{t_1-1}:K\to U_{n-1}$. Observe that $f^{-1}(K)\subset U_n$. This follows 
from the fact that the orbit of $f(u_n)$ never hits $U_{n-1}$, which was 
observed in section 2. Let $K'\subset K$
be such that $f^{t_1-1}: K'\to (-\frac34 \cdot u_{n-1}, 
                                 \frac34\cdot u_{n-1})$ is monotone and onto.
This map 
has uniform bounded distortion because it has a monotone extension up to
$U_{n-1}$. Let $K''=f^{-1}(K')\subset U_n$. The derivative  
$|Df^{t_1-1}_{|S_1}|$ can be estimated by
$$
|Df^{t_1-1}_{|S_1}|\ge C\frac{|f^{t_1}(K'')|}{|f(K'')|}
                        \ge C \frac{u_{n-1}}{u_n^\alpha}.
$$
\hfill\hfill\qed $\,\,$ (Lemma 3.5) 
\enddemo

\demo{Proof of Proposition 3.1} Assume first that $s=1$. 
This is an application of the previous Lemma 3.5. If 
$|R_n(U_{n+1})|\ge \frac{1}{10} u_n$ then
$$
|Df^{t_1}(f(x))|\ge C\cdot \frac{u_n}{u_{n+1}^\alpha}\cdot u_{n+1}^{\alpha-1}=
   C\cdot \frac{u_n}{u_{n+1}}\ge C\cdot \rho_n,
$$
where we used that $R_n(x)\notin U_{n+1}$.

\flushpar
In the other case when $|R_n(U_{n+1})|< \frac{1}{10} u_n$, we have 
$$
|Df^{t_1}(f(x))|\ge C\cdot \frac{u_{n-1}}{u_n^\alpha}\cdot u_n^{\alpha-1}
             =   C\cdot \frac{u_{n-1}}{u_{n}}\ge C\cdot \rho_n,
$$
where we used that in this case $f^{t_1-1}(x)\in R_n(U_{n+1})$
which is close to the boundary of $U_{n}$. The case with $s=1$ is finished. 

\bigskip

\flushpar
Now assume that $s\ge 2$. The proof will be split in two 
cases. Let $n\ge n_0\ge 1$ be big enough such that Lemma 3.3 and 3.4 can be 
applied. Let $K<\infty$ be the constant from Lemma 3.3 and $\theta<1$
the constant from Lemma 3.4.
 
\demo{Case I ( $f^T(0)=R_n^s(0)\in (-\frac34\cdot u_n,\frac34\cdot u_n)$)} Let 
$H\subset S_s$ be such that
$f^{T-1}=f^{t_s-1}:H\to (-\frac34 \cdot u_n, \frac34 \cdot u_n)$ is onto. 
This map has a monotone extension up to $U_n$. Hence it has a uniformly 
bounded distortion,
$$
\aligned
|Df^T(f(x))|&\ge C\cdot \frac{|f^{T-1}(H)|}{|H|}\cdot u_{n+1}^{\alpha-1}\\
            &\ge C\cdot \frac{u_n}{|S_s|}\cdot u_{n+1}^{\alpha-1}\\
            &\ge C\cdot \frac{u_n}
                             {\theta^{s-1}\cdot u_{n+1}^\alpha}
                  \cdot u_{n+1}^{\alpha-1}\\
            &\ge C\cdot \rho_n\cdot \theta^{-(s-1)}.
\endaligned
$$
\enddemo

\demo{Case II ( $f^T(0)=R_n^s(0)\notin (-\frac34\cdot u_n,\frac34\cdot u_n)$)} 
If the 
distortion of $f^{T-1}:S_s\to U_n$ is bounded by $K$ then we can give the same
argument as in case I:
$$
\aligned
|Df^T(f(x))|&\ge C\cdot \frac{|U_n|}{|S_s|}\cdot u_{n+1}^{\alpha-1}\\
            &\ge C\cdot \frac{u_n}
                             {\theta^{s-1}\cdot u_{n+1}^\alpha}
                  \cdot u_{n+1}^{\alpha-1}\\
            &\ge C\cdot \rho_n\cdot \theta^{-(s-1)}.
\endaligned
$$
Now let us assume that this distortion is bigger than $K$. Apply 
Lemma 3.3, which states $I_s\subset (-\frac34\cdot u_n,\frac34\cdot u_n)$. Then
$$
\aligned
|Df^{T}(f(0))|&=|Df^{t_{s-1}}(f(0))|\cdot |Df^{T-t_{s-1}}(f^{t_{s-1}}(f(0))|\\
              &\ge C\cdot \rho_n\cdot \theta^{-(s-2)}\cdot
                   |Df^{T-t_{s-1}}(f^{t_{s-1}}(f(0))|.
\endaligned
$$
For $s-1\ge 1$ we get this estimate from case I: 
$R_n^{s-1}(0)\in I_s\subset (-\frac34 \cdot u_n,\frac34\cdot u_n)$. 
When $s-1=1$ it follows from the Proof of Proposition 3.1 for $s=1$. 

\flushpar
The last factor can be
estimated by using the fact that $f^{T-t_{s-1}-1}:f(I_s)\to U_n$ has a monotone
extension up to $U_{n-1}$, see Lemma 2.2. It is essentially linear and its
derivative can be estimated
$$
|Df^{T-t_{s-1}-1}|_{f(I_s)}|\ge C\cdot \frac{|U_n|}{|f(I_s)|}
                            \ge C\frac{|U_n|}{\epsilon\cdot |f(U_n)|},
$$
where $\epsilon>0$ is given by Lemma 2.3. By taking $n_0\ge 1$ big enough we
can assume that $\epsilon>0$ is arbitrarily small.

\flushpar
Observe
that $|Df(f^{T-1}(f(0))|\ge C\cdot u_{n}^{\alpha-1}$ which follows from the 
fact that $|f^{T-1}(f(0))|\ge \frac34\cdot u_n$. 
We can finish the estimate for $|Df^{T}(f(0))|$ by observing that
$$
\aligned
|Df^{T-t_{s-1}}(f^{t_{s-1}}(f(0))|
&= 
|Df^{T-t_{s-1}-1}(f^{t_{s-1}}(f(0))|\cdot |Df(f^{T-1}(f(0))|\\
&\ge C\cdot \frac{|U_n|}{\epsilon\cdot |f(U_n)|}\cdot u_n^{\alpha-1}\\
&\ge C\cdot \frac{u_n}{\epsilon \cdot u_n^\alpha}\cdot u_n^{\alpha-1}
\ge C\cdot \frac{1}{\epsilon}\ge \frac{1}{\theta}.
\endaligned
$$	 
\hfill\hfill\qed $\,\,$ (Proposition 3.1)
\enddemo
\enddemo